\section{从河西的文书到洛阳的禁门：东汉的中程压力}
\label{sec:eastern-han-second-century-depth}

东汉复兴后最容易被省略的，是一段看似没有决定性覆亡的长时间。雒阳的朝廷仍在
任命官员，河西仍在转送文书，太学仍在讨论经义，边郡仍在筹粮。正因为这些事情
没有一次性地断裂，它们才构成东汉如何变重、又如何在不同地点以不同速度失去
共同节奏的过程。读这段历史不能只从后来的宦官、外戚或党锢倒推原因；应当回到
每一条路、每一批粮、每一次继承所要求的中介，看看谁能承担、谁又开始从中取得
无法被中枢轻易收回的资源。

\subsection{一项都护任命不能画出疆界}

永平年间重进西域，先从河西的西口开始。伊吾卢、鄯善和车师不在一条由洛阳直达
的线末端：军队要有马和粮，使者要有向导，绿洲王廷要权衡接待汉使会不会立即
招来另一方的报复。窦固的军队和班超的使行让东汉重新试探这些节点；七五年
陈睦遇害、若干据点撤回，则提醒人们一处驻军若无法持续取得补给，便会把原本
分散的风险集中到一座孤城。

九一年班超受命为都护，不是东汉给西域盖上一层均质的行政颜色。都护能将使节、
军援与奏报汇聚到一个职名，却仍要依赖疏勒、于阗、龟兹等绿洲的选择，以及河西
驿传和仓储的连续工作。甘英西行至安息属地未能抵达“大秦”，最能显示这种限制：
帝国可以把人送到遥远的陆路节点，不能使海路、中介商人和沿岸政权为自己的
知识和外交让路。传世史书把这些行动组织成汉廷的功绩，因而兵数、路线和各城
的服从程度都不能写得过满；但道路必须经由地方许可才可通行，是这些材料共同
显示的条件。

\begin{event}{\eventanchor{EH-0075-CHEN-MU-CRISIS}{车师告急，陈睦遇害}}{永平十八年（七五年）}{东汉 · 车师、伊吾与河西西口}{}
警报抵达河西时，朝廷面对的不是一张可以用“进”或“退”二字解决的地图。陈睦
的都护机构靠伊吾、车师和沿线驻点传递消息、储备粮马；北匈奴的压力使这些原本
分散的节点同时变成危险。继续据守意味着内地、凉州与河西都要为一条远距补给线
支付人力，撤回则意味着绿洲首领、商旅和留守者要重新判断谁还能提供保护。

《后汉书·明帝纪》《耿恭传》和《西域传》可以交叉确认危机的顺序，却都从东汉
军政的视角安排车师、北匈奴和各城。它们不能使我们替每座城规定一个共同意图。
这次撤出所解决的是一支无法持续供养的驻军的即时危险，留下的代价却是下一次
东汉返回时仍要重新取得道路、信任和向导。班超后来重建联系，并没有消除七五年
的损失，反而证明边地网络不能由一次胜败永久取得。
\end{event}

\subsection{北边不是“平定”之后的空白}

南匈奴日逐王比归附，为雒阳提供了可结盟的政治伙伴，也带来安置首领、部众和
牧地的长期事务。“内附”是一种汉廷文书的分类，不应取代不同部众的政治关系。
同一时期，羌地的冲突又把凉州、并州的军粮、流民和安置推入朝廷账目。战报能
较连续地保存将领出征和“平定”，很少留下完整的转粮清单或逃户名册；这种不
对称不是读者可以忽略的空白，而是理解中央为何不断以临时办法处理边防的线索。

一〇七年以后，陇西、金城及其周边反复告急。地方人群面对的并非一支可用单一
族名概括的力量，而是牧地、河谷、城塞、征发和保护者之间不断改变的关系。郡县
为了保城与输送而加重动员，动员又会让原已紧张的土地和人力更难维持。段颎在
一六七至一六九年的军事推进能够改变部分通道，不能证明冲突条件自此消失；战
绩所说的斩获、迁徙和“羌”的名称，也不能换成内部一致的群体或精确的社会统计。

\begin{event}{\eventanchor{EH-0107-LIANG-QIANG-WAR}{羌乱再起，凉州与并州的账不断加长}}{永初元年（一〇七年）}{东汉 · 陇西、金城、凉州与并州}{}
边郡的告急来到雒阳时，首先要变成粮、马、任命和行军的决定。谁沿河谷送粮，
谁守住城塞，谁因避战而离开原有的土地，往往不会完整进入帝纪；但这些选择决定
一支军队能否到达、一次“平定”能否维持。朝廷能够发出征调，不能同时看见每一
处因征调而失去劳力的家庭，因而军费常以文书上看不清的迁徙和债务留在地方。

《后汉书·安帝纪》《西羌传》及相关列传保存了军政处置的先后。它们不足以量化
部众、死伤或每一郡的负担，却不妨碍我们确认一个中期后果：西北边防不再只是
中央可按年拨付的边缘项目，它逐渐培养了熟悉兵源、道路与自筹资源的地方军政
网络。到灵帝朝凉州的危机重新爆发时，这些网络已经能在宫廷危机中被不同人物
调用。
\end{event}

\subsection{继承为何反复打开同一扇门}

章帝去世后，和帝年幼，窦太后临朝。幼主并不等于朝廷没有皇帝；它意味着诏令
须经过太后、近臣、尚书与守卫才能抵达官署。窦宪的北伐把边地军功和宫门内的
保护结合起来，九二年和帝与郑众诛窦氏又将这种结合反向改写。结果不是外戚从
政治中被永久排除，而是宦官掌握了守门、传诏和宫廷安全这一更可见的通道。每次
危机都解决“谁能保护皇帝”的即时问题，也令下一次危机中谁能接近皇帝变得更有
价值。

安帝死于一二五年后，阎后先立北乡侯，孙程等宦官旋即迎立刘保。传记文本对候选
人、谋划和忠奸有浓厚的立场，无法让我们重演全部宫廷选择；继统在短期内两度
改变却是可以核对的。对州郡来说，继承之所以重要，不是他们都参与了宫变，而是
任命、赏赐、赦令与文书的有效性必须重新确认。中央若反复依靠一批能控制禁门
的人完成确认，这批人就会成为官僚、外戚和地方荐举者不得不计算的中介。

\begin{event}{\eventanchor{EH-0125-AN-DI-SUCCESSION}{安帝死后，两次拥立穿过宫门}}{延光四年（一二五年）}{东汉 · 雒阳}{}
一位皇帝死去，并不会自动指出下一道诏令由谁签发。阎后先立北乡侯，孙程等人
又迎立刘保；这几天或几周中的警卫、印玺、近臣和消息，决定哪一种继承能够走到
尚书和州郡。史书所列的密谋和评价需要保留距离，但它们至少让我们看见，继统
已不能只凭宗谱完成，宫门的物理控制与文书的传递同样构成皇位的一部分。

这次处置让宦官的政治入口更深地进入制度日常，并没有让外戚、士人或地方网络
停止竞争。后来的梁氏、窦氏以及党锢危机，都是在这个条件下发生：朝廷需要有人
传递命令和保护幼主，却没有一条能让所有中介持续接受检验的安全路径。
\end{event}

\subsection{石经、察举与被截断的评价}

白虎观讨论和熹平石经，展示的不是一种已经覆盖天下的儒学国家，而是朝廷如何
把经义、书写和太学空间做成可见的公共权威。石经残片能够证明文字曾被刻立，
不能告诉我们每个士人如何读它；《白虎通》的成书也不能被当成现场速记。它们
仍说明，师承和知识可以成为接近官署的路径。察举需要地方守令的评价，士人网络
需要师友、家族和乡里声望；这些路径为中枢提供人才和信息，也可能在宫廷眼中
成为不受控制的判断中心。

梁冀被诛后，桓帝借宦官力量取得摆脱外戚的机会，新的守门者同样进入人事与赏赐
的枢纽。一六六年和一六九年的党锢，把一部分太学生、官僚与地方名士称为“党”，
但名单、地域和组织程度从来不整齐。禁令切断的不是一个拥有统一纲领的政党，
而是许多能够把地方评价送进京师的关系。短期内它压低了可见的挑战，长期却令
地方社会更难把经验安全地交给中枢，继而更依赖宗族、宾客、州郡和私人武装来
寻找保护。

\begin{sourcenote}[中期东汉的材料怎样改变读法]
《后汉书》帝纪、列传和志书将继承、边战、石经和党锢编进朝廷年序；燕然铭、
石经残片、简牍、墓葬与地方研究提供不同尺度的校验。它们不能给出连续的军粮
账、流民名册或完整社会网络。故本节不把“内附”“羌”“党人”当作统一行动者，
也不将纪功、诏令和碑刻推成所有地方的实际反应。
\end{sourcenote}
